%%=============================================================================
%% Samenvatting
%%=============================================================================

% TODO: De "abstract" of samenvatting is een kernachtige (~ 1 blz. voor een
% thesis) synthese van het document.
%
% Een goede abstract biedt een kernachtig antwoord op volgende vragen:
%
% 1. Waarover gaat de bachelorproef?
% 2. Waarom heb je er over geschreven?
% 3. Hoe heb je het onderzoek uitgevoerd?
% 4. Wat waren de resultaten? Wat blijkt uit je onderzoek?
% 5. Wat betekenen je resultaten? Wat is de relevantie voor het werkveld?
%
% Daarom bestaat een abstract uit volgende componenten:
%
% - inleiding + kaderen thema
% - probleemstelling
% - (centrale) onderzoeksvraag
% - onderzoeksdoelstelling
% - methodologie
% - resultaten (beperk tot de belangrijkste, relevant voor de onderzoeksvraag)
% - conclusies, aanbevelingen, beperkingen
%
% LET OP! Een samenvatting is GEEN voorwoord!

%%---------- Nederlandse samenvatting -----------------------------------------
%
% TODO: Als je je bachelorproef in het Engels schrijft, moet je eerst een
% Nederlandse samenvatting invoegen. Haal daarvoor onderstaande code uit
% commentaar.
% Wie zijn bachelorproef in het Nederlands schrijft, kan dit negeren, de inhoud
% wordt niet in het document ingevoegd.

\IfLanguageName{english}{%
\selectlanguage{dutch}
\chapter*{Samenvatting}

\selectlanguage{english}
}{}

%%---------- Samenvatting -----------------------------------------------------
% De samenvatting in de hoofdtaal van het document

\chapter*{\IfLanguageName{dutch}{Samenvatting}{Abstract}}



Deze bachelorproef onderzoekt in welke mate zes large language models betrouwbare en technisch correcte Python code genereren voor representatieve data analyse en ETL taken binnen de context van studenten Toegepaste Informatica aan HOGENT. Het onderzoek vertrekt vanuit de toenemende inzet van generatieve AI in programmeeropdrachten, waarbij studenten code laten genereren voor data verwerking, transformatie en automatisatie, maar waarbij niet altijd duidelijk is hoe correct, veilig en onderhoudbaar die output werkelijk is. De centrale onderzoeksvraag luidt daarom: in welke mate verschilt de technische correctheid en betrouwbaarheid van Python code die wordt gegenereerd door verschillende LLM's voor representatieve data analyse en ETL taken binnen de context van HOGENT studenten Toegepaste Informatica, en welke praktische richtlijnen kunnen studenten helpen om foutieve of suboptimale output te herkennen en te verhelpen?

Om deze vraag te beantwoorden werd een vergelijkende benchmark opgezet met zes modellen, namelijk Deepseek Coder V2, Gemma 4, llama 3, Mistral, Phi4-mini en Qwen 2.5 coder. Elk model werd getest op 35 representatieve taken voor data analyse en ETL, met drie samples per taak, wat resulteerde in 630 individuele code-evaluaties. De gegenereerde code werd volledig geautomatiseerd geëvalueerd op acht dimensies: syntaxis, style, beveiliging, uitvoerbaarheid, performantie, schaalbaarheid, onderhoudbaarheid (Radon) en functionele correctheid. De evaluatie combineerde automatische checks via AST parsing, Ruff, Bandit, runtime tests, performance profilers, Radon metrieken en functionele testen.

De resultaten tonen duidelijke verschillen tussen de modellen op het vlak van codekwaliteit, maar een opvallend consistent patroon op het vlak van functionele correctheid. De algemene rangschikking op basis van de gemiddelde kwaliteitsscore (exclusief functionele correctheid) is: deepseek-coder-v2 (90,3 procent), qwen2.5-coder (90,2 procent), gemma4 (86,5 procent), phi4-mini (85,4 procent), llama3 (81,7 procent) en mistral (79,6 procent). Opvallend is dat alle modellen uitstekende scores behaalden op syntaxis (95-100 procent), beveiliging (100 procent) en onderhoudbaarheid (boven 95 procent), maar dat de functionele correctheid sterk bleef steken tussen 29,5 procent en 58,1 procent. Gemma 4 behaalde de hoogste functionele score (58,1 procent) maar tegelijk de laagste style-compliance (54,3 procent). Qwen 2.5 coder was het enige model met perfecte style (100 procent), maar bleef steken op 57,1 procent functionele correctheid. DeepSeek Coder V2, de leider in algemene kwaliteit, behaalde eveneens 57,1 procent functionele correctheid. Llama 3, ondanks een degelijke overall ranking, behaalde slechts 41,0 procent functionele correctheid en 62,9 procent uitvoerbaarheid. Mistral toonde het grootste contrast: een lagere overall score (79,6 procent) maar eveneens de laagste functionele (29,5 procent) en execution (51,4 procent) scores.

De analyse toont aan dat functionele correctheid systematisch het grootste knelpunt vormt. De modellen genereerden regelmatig code die wel syntactisch correct en uitvoerbaar was, maar inhoudelijk toch foutieve resultaten opleverde of slechts een deel van de vereiste logica implementeerde. Uit de opsplitsing per moeilijkheidsniveau blijkt dat de functionele correctheid voor eenvoudige taken varieerde van 44,4 procent tot 75,0 procent, maar voor gemiddelde taken daalde naar 20,0 procent tot 55,0 procent. Alle zes modellen presteerden lager op gemiddelde dan op eenvoudige taken, wat bevestigt dat taken waarbij meerdere stappen gecombineerd moeten worden — zoals inlezen, transformeren, valideren en wegschrijven — systematisch vaker tot fouten leiden. De beveiligingsanalyse (Bandit) leverde een 100 procent score op voor alle modellen, wat te verklaren is doordat de ETL-taken in deze benchmark doorgaans geen beveiligingsgevoelige operaties bevatten zoals netwerkverkeer, authenticatie of cryptografie.

De meerwaarde van dit onderzoek ligt in de combinatie van een domeinspecifieke vergelijking van LLM's en een praktisch bruikbaar, volledig geautomatiseerd evaluatiekader voor studenten en docenten Toegepaste Informatica. De bevindingen tonen aan dat generatieve AI vandaag nuttig kan zijn als ondersteunende programmeerhulp, maar niet zonder menselijke controle mag worden gebruikt in data analyse en ETL contexten. Zelfs de best presterende modellen halen geen volledig overtuigend niveau, met name door de systematische zwakte op functionele correctheid. Op basis van deze bachelorproef kunnen studenten gerichter leren beoordelen wanneer LLM gegenereerde code bruikbaar is, waar de grootste risico's zitten en welke controles noodzakelijk zijn om fouten te detecteren en te verbeteren.
